\section{Теоретические сведения}
Подробное обсуждение методов поиска минимума функции в этой работе опущено,
но будут приведены основные определения и термины одномерной оптимизации, а также 
некоторые замечания по реализации алгоритмов.

\begin{definition}[Унимодальная функция]
Функция $f \colon X \to \mathbb{R}$ называется унимодальной на
интервале $X$, если существует единственная точка минимума
$x^* \in X$ и функция убывает слева от $x^*$ и возрастает справа от него.
\end{definition}


\begin{equation}
\exists x^* \in X:
\quad
\left\{
\begin{aligned}
    f(x_1) &> f(x_2),
    && \text{если } x_1 < x_2 < x^*,
    && \forall x_1, x_2 \in X; \\[0.3em]
    f(x_1) &< f(x_2),
    && \text{если } x^* < x_1 < x_2,
    && \forall x_1, x_2 \in X.
\end{aligned}
\right.
\end{equation}


Все строго выпуклые функции являются унимодальными, но не все унимодальные функции являются выпуклыми.
Некоторые разрывные функции могут обладать свойством унимодальности. 

Функции многих экстремумов носят название мультимодальных (полимодальных).
Реализация методов поиска минимума для мультимодальных функций требует дополнительных
условий, например, знания числа экстремумов на интервале или константы Липшица для функции, 
что не всегда возможно на практике. Поэтому алгоритмы поиска глобального
минимума требуют альтернативных подходов по сравнению с локальными методами.

\newpage

\begin{definition}[Липшицева функция]
Функция $f(x)$ называется липшицевой на интервале $[a,b]$, если  $\exists M>0$ такая, что:
\end{definition}

\begin{equation}
|f(x) - f(x')| \leq M|x - x'|, \qquad \forall x, x' \in X
\end{equation}

В силу того что, чтобы узнать константу Липшица $M$ точно, необходимо решать задачу 
глобальной максимизации аналитической производной функции 
(либо давать ей грубую верхнюю оценку, что тоже исключает возможность автоматизации алгоритма),
на практике используют приближённые значения константы Липшица. В работе реализован адаптивный
способ уточнения константы Липшица по уже вычисленным значениям функции, таким образом снижается 
число вычислений функции и повышается точность оценки константы Липшица с каждой итерацией метода.

Также стоит упомянуть об одной особенности действия метода парабол, который использует аппроксимацию
функции параболой через три точки, а именно поиск подходящей тройки точек для построения параболы.
Было бы неэффективно, некорректно и ресурсоемко для каждой функции вручную подбирать точки
исходя из ее графика, поэтому в работе реализован автоматический способ выбора тройки точек для
построения параболы. Вычисления функции во время этого алгоритма также учтены в общем числе вычислений
функции, что позволяет дать более объективную оценку эффективности сверхлинейного метода парабол.

\begin{definition}[Отрезок локализации]
Отрезком локализации называется отрезок
\([a,b] \subset X\), содержащий точку минимума \(x^*\):
\[
x^* \in [a,b].
\]
\end{definition}

В процессе работы методов отрезок последовательно сужается, обеспечивая уточнение
положения минимума. Бесконечношаговые алгоритмы должны иметь условия остановки,
которые обеспечивают требуемую точность координаты минимума. 
На практике часто используются следующие условия остановки:
\begin{gather}
    \|x^{k+1} - x^k\| \leq \varepsilon_1; \\
    \|f(x^{k+1}) - f(x^k)\| \leq \varepsilon_2; \\
    \|f'(x^{k+1})\| \leq \varepsilon_3.
\end{gather}
Вообще говоря выполнение вышеуказанных условий не гарантирует нахождение
глобального минимума, но в случаеунимодальной функции, 
если алгоритм корректно реализован, выполнение этих условий
гарантирует достижения необходимой точности решения задачи.
